\textbf{Phần 1. Trắc nghiệm lựa chọn}
\setcounterpageref{ex}{0}

\begin{ex} %[5D4V24]
Giá trị của biểu thức $12{,}8 \times 4{,}5 + 12{,}8 \times 5{,}5 - 28$ là
\choice
{$90$}
{\True $100$}
{$110$}
{$128$}
\loigiai{
Ta nhóm thừa số chung $12{,}8$:
$$12{,}8 \times (4{,}5 + 5{,}5) - 28 = 12{,}8 \times 10 - 28 = 128 - 28 = 100.$$
\textbf{Đáp án của câu hỏi là B}
}
\end{ex}

\begin{ex} %[5D4V24]
Tìm $x$, biết: $(x + 3{,}5) \times 1{,}2 = 9$.
\choice
{$3$}
{\True $4$}
{$4{,}5$}
{$5$}
\loigiai{
Ta có:
\begin{align*}
x + 3{,}5 &= 9 : 1{,}2\\
x + 3{,}5 &= 7{,}5\\
x &= 7{,}5 - 3{,}5 = 4
\end{align*}
\textbf{Đáp án của câu hỏi là B}
}
\end{ex}

\begin{ex} %[5D4V24]
Một mảnh vườn hình chữ nhật có chu vi $48\text{ m}$, chiều dài hơn chiều rộng $4{,}8\text{ m}$. Diện tích của mảnh vườn đó là
\choice
{$144\text{ m}^2$}
{\True $138{,}24\text{ m}^2$}
{$136{,}44\text{ m}^2$}
{$142{,}5\text{ m}^2$}
\loigiai{
Nửa chu vi của mảnh vườn là:
$48 : 2 = 24\text{ (m)}$.\\
Chiều rộng mảnh vườn là:
$(24 - 4{,}8) : 2 = 9{,}6\text{ (m)}$.\\
Chiều dài mảnh vườn là:
$24 - 9{,}6 = 14{,}4\text{ (m)}$.\\
Diện tích mảnh vườn là:
$14{,}4 \times 9{,}6 = 138{,}24\text{ (m}^2\text{)}$.\\
\textbf{Đáp án của câu hỏi là B}
}
\end{ex}

\begin{ex} %[5D4V24]
Một người đi xe đạp, trong 3 giờ đầu mỗi giờ đi được $12{,}5\text{ km}$; trong 2 giờ sau mỗi giờ đi được $11{,}25\text{ km}$. Hỏi trung bình mỗi giờ người đó đi được bao nhiêu ki-lô-mét?
\choice
{$11{,}875\text{ km}$}
{\True $12\text{ km}$}
{$12{,}25\text{ km}$}
{$12{,}5\text{ km}$}
\loigiai{
Quãng đường đi trong 3 giờ đầu là:
$12{,}5 \times 3 = 37{,}5\text{ (km)}$.\\
Quãng đường đi trong 2 giờ sau là:
$11{,}25 \times 2 = 22{,}5\text{ (km)}$.\\
Tổng quãng đường đi được là:
$37{,}5 + 22{,}5 = 60\text{ (km)}$.\\
Tổng thời gian đi là:
$3 + 2 = 5\text{ (giờ)}$.\\
Trung bình mỗi giờ người đó đi được:
$60 : 5 = 12\text{ (km)}$.\\
\textbf{Đáp án của câu hỏi là B}
}
\end{ex}

\begin{ex} %[5D4V24]
Mỗi chai chứa $0{,}75\text{ l}$ sữa tươi. Mỗi lít sữa tươi cân nặng $1{,}08\text{ kg}$. Biết vỏ chai rỗng nặng $0{,}19\text{ kg}$. Hỏi $12$ chai sữa như thế cân nặng tất cả bao nhiêu ki-lô-gam?
\choice
{$10{,}8\text{ kg}$}
{\True $12\text{ kg}$}
{$12{,}4\text{ kg}$}
{$11{,}6\text{ kg}$}
\loigiai{
Lượng sữa trong một chai cân nặng số ki-lô-gam là:
$0{,}75 \times 1{,}08 = 0{,}81\text{ (kg)}$.\\
Một chai sữa (gồm cả vỏ) cân nặng số ki-lô-gam là:
$0{,}81 + 0{,}19 = 1\text{ (kg)}$.\\
$12$ chai sữa như thế cân nặng số ki-lô-gam là:
$1 \times 12 = 12\text{ (kg)}$.\\
\textbf{Đáp án của câu hỏi là B}
}
\end{ex}

\begin{ex} %[5D4V24]
Khi cộng một số tự nhiên với một số thập phân có một chữ số ở phần thập phân, một bạn học sinh đã quên viết dấu phẩy của số thập phân nên thực hiện cộng như hai số tự nhiên và được kết quả là $492$. Tổng đúng của phép tính là $141{,}9$. Số tự nhiên cần tìm là
\choice
{$38{,}9$}
{$102$}
{\True $103$}
{$104$}
\loigiai{
Khi bỏ quên dấu phẩy ở số thập phân có một chữ số thập phân, số đó đã gấp lên $10$ lần.\\
Do đó tổng tăng thêm một lượng bằng $9$ lần số thập phân đó.\\
Số thập phân đó là:
$(492 - 141{,}9) : 9 = 350{,}1 : 9 = 38{,}9$.\\
Số tự nhiên cần tìm là:
$141{,}9 - 38{,}9 = 103$.\\
\textbf{Đáp án của câu hỏi là C}
}
\end{ex}

\begin{ex} %[5D4V24]
Tính giá trị của biểu thức $(25{,}4 - 13{,}8) \times 2{,}5 + 1{,}25 \times 8$:
\choice
{$38$}
{\True $39$}
{$40$}
{$41$}
\loigiai{
Ta có:
$$(25{,}4 - 13{,}8) \times 2{,}5 + 1{,}25 \times 8 = 11{,}6 \times 2{,}5 + 10 = 29 + 10 = 39.$$
\textbf{Đáp án của câu hỏi là B}
}
\end{ex}

\begin{ex} %[5D4V24]
Người ta quét vôi xung quanh bốn bức tường bên trong một phòng học hình chữ nhật dài $8\text{ m}$, rộng $6\text{ m}$ và cao $3{,}5\text{ m}$. Tổng diện tích các cửa sổ và cửa ra vào là $9\text{ m}^2$. Diện tích cần quét vôi là
\choice
{$98\text{ m}^2$}
{\True $89\text{ m}^2$}
{$95\text{ m}^2$}
{$80\text{ m}^2$}
\loigiai{
Chu vi đáy của phòng học là:
$(8 + 6) \times 2 = 28\text{ (m)}$.\\
Diện tích xung quanh bốn bức tường là:
$28 \times 3{,}5 = 98\text{ (m}^2\text{)}$.\\
Diện tích cần quét vôi là:
$98 - 9 = 89\text{ (m}^2\text{)}$.\\
\textbf{Đáp án của câu hỏi là B}
}
\end{ex}

\begin{ex} %[5D4V24]
Mua 4 hộp sữa tươi và 3 gói bánh hết $116\,000$ đồng. Mua 2 hộp sữa tươi và 3 gói bánh cùng loại hết $84\,000$ đồng. Giá tiền một hộp sữa tươi là
\choice
{$14\,000\text{ đồng}$}
{\True $16\,000\text{ đồng}$}
{$18\,000\text{ đồng}$}
{$20\,000\text{ đồng}$}
\loigiai{
Hai hộp sữa tươi có giá tiền là:
$116\,000 - 84\,000 = 32\,000\text{ (đồng)}$.\\
Giá tiền của một hộp sữa tươi là:
$32\,000 : 2 = 16\,000\text{ (đồng)}$.\\
\textbf{Đáp án của câu hỏi là B}
}
\end{ex}

\begin{ex} %[5D4V24]
Tìm giá trị của $x$, biết: $x \times 0{,}75 + x \times 0{,}25 = 14{,}6 \times 5$.
\choice
{$70$}
{$72$}
{\True $73$}
{$75$}
\loigiai{
Ta có:
\begin{align*}
x \times (0{,}75 + 0{,}25) &= 73\\
x \times 1 &= 73\\
x &= 73
\end{align*}
\textbf{Đáp án của câu hỏi là C}
}
\end{ex}

\begin{ex} %[5D4V24]
Bác nông dân thu hoạch được $2{,}4$ tấn lúa tươi. Sau khi phơi khô, khối lượng lúa bị hao hụt mất $0{,}1$ khối lượng lúa tươi. Bác nông dân thu được bao nhiêu tạ lúa khô?
\choice
{$21\text{ tạ}$}
{\True $21{,}6\text{ tạ}$}
{$22\text{ tạ}$}
{$22{,}4\text{ tạ}$}
\loigiai{
Khối lượng lúa hao hụt là:
$2{,}4 \times 0{,}1 = 0{,}24\text{ (tấn)}$.\\
Khối lượng lúa khô thu được là:
$2{,}4 - 0{,}24 = 2{,}16\text{ (tấn)}$.\\
Đổi: $2{,}16\text{ tấn} = 21{,}6\text{ tạ}$.\\
\textbf{Đáp án của câu hỏi là B}
}
\end{ex}

\begin{ex} %[5D4V24]
Tìm số dư của phép chia $43{,}19 : 2{,}8$ khi lấy phần thập phân của thương đến một chữ số.
\choice
{$0{,}7$}
{\True $0{,}07$}
{$0{,}007$}
{$7$}
\loigiai{
Thực hiện phép chia: $43{,}19 : 2{,}8 = 15{,}4$ (lấy một chữ số ở phần thập phân).\\
Ta có: $15{,}4 \times 2{,}8 = 43{,}12$.\\
Số dư là: $43{,}19 - 43{,}12 = 0{,}07$.\\
\textbf{Đáp án của câu hỏi là B}
}
\end{ex}

\textbf{Phần 2. Trắc nghiệm đúng sai}
\setcounterpageref{ex}{0}

\begin{ex} %[5D4V24]
Cho hai biểu thức $P = (4{,}5 + 2{,}5) \times 3{,}6$ và $Q = 90 : 2{,}5 - 10{,}8$.
\choiceTF
{\True Giá trị của biểu thức $P$ bằng $25{,}2$}
{\True Giá trị của biểu thức $Q$ bằng $25{,}2$}
{Giá trị của biểu thức $P$ lớn hơn giá trị của biểu thức $Q$}
{\True Tích $P \times 10$ bằng $252$}
\loigiai{
\begin{itemchoice}
\item Đúng. $P = 7 \times 3{,}6 = 25{,}2$.
\item Đúng. $Q = 36 - 10{,}8 = 25{,}2$.
\item Sai. Vì $P = Q = 25{,}2$.
\item Đúng. $25{,}2 \times 10 = 252$.
\end{itemchoice}
}
\end{ex}

\begin{ex} %[5D4V24]
Một bể kính nuôi cá dạng hình hộp chữ nhật có lòng trong dài $1{,}2\text{ m}$, rộng $0{,}8\text{ m}$ và cao $0{,}6\text{ m}$.
\choiceTF
{\True Thể tích của lòng bể kính là $0{,}576\text{ m}^3$}
{\True Thể tích đó tương đương với $576\text{ l}$ nước}
{\True Mức nước hiện tại trong bể cao bằng $\dfrac{3}{4}$ chiều cao bể, tương ứng với chiều cao mực nước là $0{,}45\text{ m}$}
{Lượng nước hiện đang có trong bể là $400\text{ l}$}
\loigiai{
\begin{itemchoice}
\item Đúng. Thể tích lòng bể là: $1{,}2 \times 0{,}8 \times 0{,}6 = 0{,}576\text{ (m}^3\text{)}$.
\item Đúng. Vì $1\text{ m}^3 = 1\,000\text{ dm}^3 = 1\,000\text{ l}$ nên $0{,}576\text{ m}^3 = 576\text{ l}$.
\item Đúng. Chiều cao mực nước là: $0{,}6 \times \dfrac{3}{4} = 0{,}45\text{ (m)}$.
\item Sai. Lượng nước hiện có trong bể là: $576 \times \dfrac{3}{4} = 432\text{ (l)} 
eq 400\text{ l}$.
\end{itemchoice}
}
\end{ex}

\textbf{Phần 3. Trắc nghiệm trả lời ngắn}
\setcounterpageref{ex}{0}

\begin{ex} %[5D4V24]
Tính giá trị của biểu thức: $0{,}4 \times 125 \times 2{,}5 \times 0{,}8$.
\shortans{100}
\loigiai{
Áp dụng tính chất giao hoán và kết hợp:
$$(0{,}4 \times 2{,}5) \times (125 \times 0{,}8) = 1 \times 100 = 100.$$
}
\end{ex}

\begin{ex} %[5D4V24]
Tính tổng các giá trị tự nhiên của $x$ thỏa mãn: $18{,}4 < 3{,}2 \times x < 24$.
\shortans{13}
\loigiai{
Ta có: $18{,}4 : 3{,}2 < x < 24 : 3{,}2 \implies 5{,}75 < x < 7{,}5$.\\
Các số tự nhiên $x$ thỏa mãn là: $6$ và $7$.\\
Tổng của chúng là: $6 + 7 = 13$.
}
\end{ex}

\begin{ex} %[5D4V24]
May 8 bộ quần áo trẻ em đồng phục hết $14\text{ m}$ vải. Hỏi với $35\text{ m}$ vải cùng loại thì may được bao nhiêu bộ quần áo đồng phục như thế?
\shortans{20}
\loigiai{
May một bộ quần áo cần số mét vải là:
$14 : 8 = 1{,}75\text{ (m)}$.\\
Số bộ quần áo may được là:
$35 : 1{,}75 = 20\text{ (bộ)}$.
}
\end{ex}

\begin{ex} %[5D4V24]
Một mảnh ruộng hình thang có đáy lớn $24\text{ m}$, đáy bé $16\text{ m}$ và chiều cao $12{,}5\text{ m}$. Diện tích của mảnh ruộng đó là bao nhiêu mét vuông?
\shortans{250}
\loigiai{
Diện tích của mảnh ruộng hình thang là:
$(24 + 16) \times 12{,}5 : 2 = 40 \times 12{,}5 : 2 = 250\text{ (m}^2\text{)}$.
}
\end{ex}

\begin{ex} %[5D4V24]
Một xe ô tô khởi hành từ tỉnh A lúc 7 giờ và đến tỉnh B lúc 10 giờ 30 phút. Trên đường đi ô tô nghỉ giải lao 30 phút. Vận tốc của xe là $48\text{ km/h}$. Tính độ dài quãng đường AB theo đơn vị ki-lô-mét.
\shortans{144}
\loigiai{
Thời gian ô tô thực tế di chuyển trên đường là:
$10\text{ giờ }30\text{ phút} - 7\text{ giờ} - 30\text{ phút} = 3\text{ giờ}$.\\
Quãng đường AB dài:
$48 \times 3 = 144\text{ (km)}$.
}
\end{ex}

\begin{ex} %[5D4V24]
Trung bình cộng của ba số bằng $25$. Số thứ nhất là $18{,}5$, số thứ hai là $26{,}5$. Tìm số thứ ba.
\shortans{30}
\loigiai{
Tổng của ba số là:
$25 \times 3 = 75$.\\
Số thứ ba là:
$75 - (18{,}5 + 26{,}5) = 75 - 45 = 30$.
}
\end{ex}
